\section{Solutions for Chapter~\ref{ch:modular}}

\paragraph{Exercise~\ref{ex:Z3-S-square}.}
\phantomsection
\label{sol:ex:Z3-S-square}
Summing over the intermediate label $(a,b)$ gives
\begin{align*}
 (S^2)_{(g,k),(h,l)}
 &=\frac19\sum_{a,b}\omega^{-ka-bg-bh-la}\\
 &=\left(\frac13\sum_a\omega^{-(k+l)a}\right)
   \left(\frac13\sum_b\omega^{-(g+h)b}\right)\\
 &=\delta_{l,-k}\delta_{h,-g}.
\end{align*}
For example $(1,0)$ is conjugate to $(2,0)$, rather than to
itself.  Therefore charge conjugation is not the identity matrix.



\paragraph{Exercise~\ref{ex:H3-S-square}.}
\phantomsection
\label{sol:ex:H3-S-square}
Let $\mathbf1_m$ be the length-$m$ all-ones column and $J_{m,n}$
the all-ones rectangular matrix. Direct arithmetic gives
\[
 AB=\frac1{36}-\frac1{52}=\frac1{117},\qquad
 A^2+B^2=\frac1{18}+\frac1{26}=\frac{11}{117}.
\]
Using $u^2=1/9$ and $v^2=1/13$, the upper-left diagonal entries
of $S^2$ are
$11/117+4/9+6/13=1$, and its off-diagonal entries are
$2/117+4/9-6/13=0$.
The row sums $C_4\mathbf1_4=-u\mathbf1_4$ and
$D_6\mathbf1_6=v\mathbf1_6$, together with
$A+B=u$ and $A-B=-v$, give the first two rows' mixed blocks:
\[
 u(A+B-u)\mathbf1_4^{\mathsf t}=0,
 \quad v(A-B+v)\mathbf1_6^{\mathsf t}=0,
 \quad v(B-A-v)\mathbf1_6^{\mathsf t}=0.
\]
The $4\times6$ block is
$uvJ_{4,6}-uvJ_{4,6}=0$. Finally,
\[
 C_4^2+2u^2J_4=I_4,\qquad
 D_6^2+2v^2J_6=I_6
\]
by the Fourier sums in the chapter. Symmetry supplies the
transposed blocks, completing the entire identity $S^2=I$.



\paragraph{Exercise~\ref{ex:Gauss13}.}
\phantomsection
\label{sol:ex:Gauss13}
The squares are $\{1,3,4,9,10,12\}$ and the nonsquares are
$\{2,5,6,7,8,11\}$.  The list $6j^2$ for $j=1,\ldots,6$
is $6,11,2,5,7,8$ modulo thirteen.  For the Gauss--Milgram
calculation, $1+\omega+\omega^2=0$ combines the four
dimension-$s$ terms to $s^2$. Since $s=1+d^2$ and
$d^4=11d^2-1$,
\[
 1+d^4+s^2=2+2d^2+2d^4=24d^2.
\]
The six residual twists contribute
$9d^2(-1-\sqrt{13})/2$, so
\[
 p^+=\frac{d^2}{2}(39-9\sqrt{13}).
\]
Using $d^2=(11+3\sqrt{13})/2$ and multiplying out gives
$p^+=(39+9\sqrt{13})/2=3s$.



\paragraph{Exercise~\ref{ex:Verlinde-d1}.}
\phantomsection
\label{sol:ex:Verlinde-d1}
For $i=1$, $[2i]=2$; for $i=4$, $[8]=5$ because $8\equiv-5$
modulo thirteen. Hence
\begin{align*}
 d_1^2&=\one+a+c_0+c_1+c_2+d_1+d_3+d_4+d_5+d_6,\\
 d_4^2&=\one+a+c_0+c_1+c_2+d_1+d_2+d_3+d_4+d_6.
\end{align*}
The absent residual objects are $d_2$ and $d_5$, respectively.
Both dimension sums equal
$1+4s+15d=9+27d=9d^2$.



\paragraph{Exercise~\ref{ex:modular-extension}.}
\phantomsection
\label{sol:ex:modular-extension}
The first two components of $Sm$ are each $A+B+2u=1$.
The $a$ component is $2u+4/3=2$; each of the three $c_j$
components is $2u-2/3=0$.  Every final component is $v-v=0$.
Thus $Sm=m$.  Each nonzero entry of $m$ lies at twist one, so
$\Theta m=m$.  Symmetry of $S$ gives
$Smm^{\mathsf t}=mm^{\mathsf t}=mm^{\mathsf t}S$, and
the analogous statement for $\Theta$ is immediate.
To build a local CFT one still needs, among other things, actual
chiral modules and vertex operators, a compatible algebra
multiplication on the extension, and the appropriate positivity
and sewing results.  A modular-invariant integer matrix contains
none of those multiplication maps by itself.



\paragraph{Exercise~\ref{ex:central-charge-scope}.}
\phantomsection
\label{sol:ex:central-charge-scope}
In the first statement, the \emph{braided chiral representation
category} $\Rep(U)$ is a Drinfeld center. Its Gauss--Milgram phase
is one, so the chiral central charge obeys
$e^{2\pi ic_U/8}=1$. The value $c_U=8$ satisfies this constraint.
In the second statement, $\HH_3$ is a fusion category of full-CFT
symmetry lines. That does not identify the theory's entire chiral
representation category with $Z(\HH_3)$. Its center organizes
symmetry sectors, which is a different assertion. Hence a proposed
nonchiral theory with $c_L=c_R=2$ has no contradiction with the
first construction: its left-right gravitational anomaly vanishes,
and the symmetry claim alone imposes no condition $c_L=0\bmod8$.
This explains compatibility; it does not establish the proposed
$c=2$ theory or an anyon-chain scaling limit.


